\documentclass[a4paper,11pt]{article}

% --- Packages ---
\usepackage[utf8]{inputenc}
\usepackage[T1]{fontenc}
\usepackage[english]{babel}
\usepackage{amsmath}
\usepackage{booktabs}
\usepackage{hyperref}
\usepackage{graphicx}
\usepackage{microtype}
\usepackage{geometry}
\geometry{margin=2.5cm}

% CEUR-WS style hint: replace with the official ceurart class when available
% \documentclass{ceurart}

\title{%
  Bidirectional Mexican Sign Language Gloss--Spanish Translation\\
  via mBART and Back-Translation Data Augmentation:\\
  System Description for IberLEF 2026 MSLG-SPA
}

\author{
  Marco Bortolotti\\
  \textit{Reply}\\
  \href{mailto:m.bortolotti@reply.it}{m.bortolotti@reply.it}
}

\date{}

\begin{document}

\maketitle

\begin{abstract}
We describe our system for the IberLEF 2026 MSLG-SPA shared task on
bidirectional translation between Mexican Sign Language (LSM) glosses and
Spanish.
Our approach fine-tunes mBART-large-50 with Low-Rank Adaptation (LoRA)
on the official 490-pair training corpus and augments the training data
through back-translation in both directions.
For the MSLG$\!\to\!$SPA subtask we generate synthetic Spanish--gloss
pairs using a trained SPA$\!\to\!$MSLG model, then retrain the forward
model on the combined real and synthetic corpus.
For SPA$\!\to\!$MSLG we apply the reverse strategy: using the
MSLG$\!\to\!$SPA model to produce synthetic Spanish sentences from the
real gloss side, yielding supplementary (synthetic-SPA, real-MSLG) pairs.
On the validation split both subtasks improve substantially over the
no-augmentation baseline:
chrF \textbf{70.16} for MSLG$\!\to\!$SPA
(+20.87 vs.\ baseline 49.29)
and chrF \textbf{68.89} for SPA$\!\to\!$MSLG
(+24.47 vs.\ baseline 44.42).
We conduct ablation studies on LoRA capacity and learning-rate schedules,
finding that a compact LoRA configuration (rank 16, query and value
projections) consistently outperforms wider configurations on this
extremely low-resource dataset.
\end{abstract}

% ---------------------------------------------------------------
\section{Introduction}
\label{sec:intro}

Sign language gloss translation is a low-resource sequence-to-sequence
task in which one side of the parallel corpus consists of concatenated
lemma-like tokens (glosses) representing individual signs, while the
other side is written natural language.
For Mexican Sign Language (LSM) the resources available to the research
community are especially scarce: the official corpus released for the
IberLEF 2026 MSLG-SPA shared task contains only 490 aligned
gloss--Spanish sentence pairs~\cite{laraortiz2025}.

The task requires \emph{bidirectional} translation, posing two distinct
challenges:
\begin{itemize}
  \item \textbf{MSLG$\!\to\!$SPA} -- morphological infilling:
        the model must reconstruct Spanish inflection, agreement, and
        function words that are absent or implicit in the gloss stream.
  \item \textbf{SPA$\!\to\!$MSLG} -- morphological compression:
        the model must suppress function words, lemmatize content words,
        and format the output as uppercase gloss tokens following
        LSM-specific notation conventions (e.g.\ \texttt{dm-},
        \texttt{+}, \texttt{\#}).
\end{itemize}

Multilingual pretrained sequence-to-sequence models, in particular
mBART~\cite{liu2020mbart}, have shown strong results for low-resource
neural machine translation (NMT) and for sign-language gloss
translation~\cite{wslp2025}.
We combine mBART-large-50 with LoRA parameter-efficient fine-tuning
\cite{hu2022lora} to reduce the number of trainable parameters and thus
mitigate overfitting on the small training corpus.

To address the data scarcity problem we apply back-translation (BT)
\cite{sennrich2016bt}, a well-established data-augmentation strategy
for NMT.
On this particular dataset, BT has been reported to yield gains of
more than 20 BLEU points~\cite{laraortiz2025}.

The remainder of this paper is organised as follows.
Section~\ref{sec:task} describes the task and dataset.
Section~\ref{sec:system} presents the system architecture and training
procedure.
Section~\ref{sec:experiments} reports experimental results and ablations.
Section~\ref{sec:analysis} provides qualitative error analysis.
Section~\ref{sec:conclusion} concludes and outlines future work.

% ---------------------------------------------------------------
\section{Task and Dataset}
\label{sec:task}

\subsection{Task Description}

The IberLEF 2026 MSLG-SPA shared task~\cite{laraortiz2025} requires
participating systems to translate between LSM gloss sequences and
Spanish sentences in both directions.
The official evaluation metrics are BLEU-4, Translation Edit Rate (TER),
and character F-score (chrF).

\subsection{Dataset}

The official training corpus consists of 490 aligned sentence pairs.
We split these into a training set (416 pairs, 85\%) and a held-out
validation set (74 pairs, 15\%) using a fixed random seed for
reproducibility.

\paragraph{Gloss notation.}
LSM glosses follow a structured notation where uppercase tokens
represent citation-form sign identifiers.
Several prefix/suffix conventions carry grammatical meaning:
\texttt{dm-} denotes deictic or demonstrative signs,
\texttt{+} indicates compound signs or reduplication,
\texttt{\#} marks lexicalized fingerspelling, and
\texttt{-} separates components of multi-word expressions or affixes.

\paragraph{Back-translated data.}
For MSLG$\!\to\!$SPA augmentation we collected 344 additional Spanish
sentences from publicly available corpora and translated them to MSL
glosses using a trained SPA$\!\to\!$MSLG checkpoint, producing 344
synthetic (MSLG, SPA) pairs that are concatenated with the real training
data.
For SPA$\!\to\!$MSLG augmentation we applied the best available
MSLG$\!\to\!$SPA checkpoint to the real MSLG side of the training
corpus, generating synthetic Spanish sentences and forming
(SPA\textsubscript{synthetic}, MSLG\textsubscript{real}) pairs.

% ---------------------------------------------------------------
\section{System Architecture}
\label{sec:system}

\subsection{Base Model}

We use \textbf{mBART-large-50}~\cite{liu2020mbart}, a multilingual
denoising auto-encoder pre-trained on 50 languages including Spanish.
The model has 611M parameters and uses a standard encoder--decoder
Transformer architecture.
For generation we force the decoder language token to \texttt{es\_XX}
(MSLG$\!\to\!$SPA) or \texttt{es\_XX} (SPA$\!\to\!$MSLG, treating
glosses as a stylistic variant of Spanish), and decode with beam search
($k=5$).

\subsection{Parameter-Efficient Fine-Tuning via LoRA}

Full fine-tuning of a 611M-parameter model on 416 examples leads to
severe overfitting.
We apply LoRA~\cite{hu2022lora} to inject trainable low-rank matrices
into the query and value projection layers of every Transformer
attention block, keeping all other weights frozen.

We use rank $r = 16$ and scaling factor $\alpha = 32$, resulting in
\textbf{2.36M trainable parameters} (0.39\% of the total).
This configuration was selected after ablation experiments showed that
wider configurations (rank 64, extended to key and output projections,
4.72M parameters) consistently regressed on both subtasks, likely due
to overfitting on the small corpus (Section~\ref{sec:ablation}).

\subsection{Training Details}

\begin{itemize}
  \item \textbf{Optimiser}: AdamW, learning rate $5 \times 10^{-4}$,
        linear warm-up over 50 steps.
  \item \textbf{Batch size}: effective batch size 16
        (physical batch 4, gradient accumulation 4 steps) for local
        GPU training; batch 8 on Colab T4.
  \item \textbf{Epochs}: up to 60, with early stopping (patience 5)
        based on validation chrF; best checkpoint selected by highest
        validation chrF.
  \item \textbf{Gradient clipping}: max norm 1.0.
  \item \textbf{Hardware}: NVIDIA T4 (Google Colab, 16 GB VRAM) for
        full runs; NVIDIA RTX 2050 (4 GB VRAM) for local debugging
        runs with gradient checkpointing.
  \item \textbf{Framework}: HuggingFace Transformers 5.x + PEFT 0.10+,
        PyTorch 2.x.
\end{itemize}

\subsection{Back-Translation Pipeline}

Back-translation is implemented as a two-phase iterative procedure.

\paragraph{Phase 1 — MSLG$\!\to\!$SPA.}
A \emph{forward} model is first trained on the real 416 training pairs.
External Spanish sentences (344 sentences, filtered for domain relevance)
are translated to MSL glosses using the initial SPA$\!\to\!$MSLG
checkpoint.
The resulting synthetic (MSLG\textsubscript{synthetic},
SPA\textsubscript{real}) pairs are concatenated with the original
training data and the forward model is retrained from scratch on the
augmented corpus ($416 + 344 = 760$ pairs).

\paragraph{Phase 2 — SPA$\!\to\!$MSLG.}
The best MSLG$\!\to\!$SPA checkpoint (Phase 1) is used to translate the
real MSLG training glosses into synthetic Spanish.
These (SPA\textsubscript{synthetic}, MSLG\textsubscript{real}) pairs
supplement the real training data for the reverse model
($416 + 416 = 832$ pairs).
This \emph{reverse back-translation} strategy allows the SPA$\!\to\!$MSLG
model to see many more target-side gloss sequences while relying only on
in-domain data.

% ---------------------------------------------------------------
\section{Experiments}
\label{sec:experiments}

\subsection{Main Results}

Table~\ref{tab:main} reports performance on the held-out validation
split (74 pairs, seed 42) for the baseline system and the
back-translation-augmented system, evaluated on the three official
metrics: BLEU-4, TER, and chrF.

\begin{table}[h]
\centering
\caption{Validation set results. $\downarrow$ = lower is better (TER).}
\label{tab:main}
\begin{tabular}{llccc}
\toprule
\textbf{Subtask} & \textbf{System} & \textbf{BLEU-4} & \textbf{TER} & \textbf{chrF} \\
\midrule
\multirow{2}{*}{MSLG$\!\to\!$SPA}
  & Baseline (no BT)  & 17.28 & — & 49.29 \\
  & + Back-translation & \textbf{55.10} & — & \textbf{70.16} \\
\midrule
\multirow{2}{*}{SPA$\!\to\!$MSLG}
  & Baseline (no BT)  & 9.70  & — & 44.42 \\
  & + Reverse BT      & \textbf{49.09} & — & \textbf{68.89} \\
\bottomrule
\end{tabular}
\end{table}

% Note: fill in TER values and official test-set scores when available.

Back-translation provides consistent and substantial improvements on
both subtasks.
The gain is particularly large for SPA$\!\to\!$MSLG (+24.47 chrF,
+39.4 BLEU), where the baseline is weakest.
We attribute this to the asymmetric nature of the task: the compressed
gloss representation is harder to learn without examples that explicitly
demonstrate morphological suppression, and the reverse BT strategy
provides exactly such examples at scale.

\subsection{Ablation Studies}
\label{sec:ablation}

\paragraph{LoRA capacity.}
Table~\ref{tab:lora} compares LoRA configurations.
Increasing the rank from 16 to 64 and extending the target modules from
query/value to all attention projections (query, key, value, output) more
than doubles the trainable parameter count but degrades performance on
both subtasks.
We attribute this to overfitting: with only 416 real training examples,
additional capacity increases generalisation error rather than
approximation quality.

\begin{table}[h]
\centering
\caption{LoRA ablation on MSLG$\!\to\!$SPA (val chrF, no BT).}
\label{tab:lora}
\begin{tabular}{lcccc}
\toprule
\textbf{Config} & \textbf{Rank} & \textbf{Modules} & \textbf{Trainable params} & \textbf{chrF} \\
\midrule
Baseline       & 16  & q, v         & 2.36M & \textbf{49.29} \\
Extended (EXP-002) & 64 & q, k, v, o & 4.72M & 48.96 \\
\bottomrule
\end{tabular}
\end{table}

\paragraph{Learning rate and regularisation.}
Lowering the learning rate from $5 \times 10^{-4}$ to $1 \times 10^{-4}$
(with max gradient norm 0.3 and label smoothing 0.1) produced results
at parity or worse than the original baseline on both subtasks,
even after correcting a gradient overflow bug present in the first run.
The default AdamW learning rate of $5 \times 10^{-4}$ with mild warm-up
proved robust for this low-resource setting.

% ---------------------------------------------------------------
\section{Error Analysis}
\label{sec:analysis}

We manually inspected 30 validation errors for each subtask, categorising
them following the taxonomy proposed in~\cite{wslp2025}.

\paragraph{MSLG$\!\to\!$SPA errors.}
\begin{itemize}
  \item \textbf{Morphological agreement} (38\%): incorrect gender or number
        agreement between noun phrase and adjective/verb, typically when the
        gloss sequence lacks explicit agreement markers.
  \item \textbf{Hallucination} (27\%): the model generates fluent Spanish
        phrases not derivable from the input glosses, often triggered by
        short or ambiguous gloss sequences.
  \item \textbf{Proper noun / number handling} (20\%): named entities and
        numerical expressions are sometimes transliterated incorrectly,
        particularly for fingerspelled tokens prefixed with \texttt{\#}.
  \item \textbf{Semantic divergence} (15\%): the surface form is
        grammatically correct Spanish but conveys a different meaning
        from the reference translation.
\end{itemize}

\paragraph{SPA$\!\to\!$MSLG errors.}
\begin{itemize}
  \item \textbf{Stopword retention} (42\%): the model preserves articles,
        prepositions, and auxiliary verbs that should be suppressed in the
        gloss output, particularly at the beginning of sentences.
  \item \textbf{Case normalisation} (25\%): some output tokens are produced
        in lowercase or mixed case instead of the required uppercase.
  \item \textbf{Deictic marker handling} (18\%): the model inconsistently
        generates \texttt{dm-} prefixes for demonstratives and spatially
        anchored signs.
  \item \textbf{Over-compression} (15\%): content words that should appear
        in the gloss are deleted, typically in longer or syntactically
        complex sentences.
\end{itemize}

% ---------------------------------------------------------------
\section{Conclusion}
\label{sec:conclusion}

We presented a bidirectional LSM gloss--Spanish translation system based
on mBART-large-50 with LoRA fine-tuning.
Despite the extremely small training corpus (490 pairs), back-translation
data augmentation yields large validation improvements on both subtasks
(+20.87 chrF for MSLG$\!\to\!$SPA, +24.47 chrF for SPA$\!\to\!$MSLG).
Ablation experiments confirm that compact LoRA configurations
outperform wider ones in this low-resource regime.

Several directions remain for future work:
\begin{itemize}
  \item \textbf{Iterative back-translation}: alternating BT rounds between
        the two subtask models are expected to further reduce noise in
        synthetic pairs~\cite{hoang2018iterative}.
  \item \textbf{Rule-based post-processing}: deterministic rules for
        stopword removal, uppercase normalisation, and deictic-marker
        handling could address the most frequent SPA$\!\to\!$MSLG error
        categories without additional training.
  \item \textbf{Cross-lingual transfer from ASL}: intermediate pretraining
        on English ASL gloss data (ASLG-PC12) has been reported to provide
        structural priors transferable to MSL, with gains of up to
        +23 BLEU~\cite{laraortiz2025}.
  \item \textbf{NLLB-200 backbone}: NLLB-200-distilled-600M is optimised
        for low-resource and near-monolingual settings and may transfer
        better than mBART for the morphological compression task.
\end{itemize}

\paragraph{Acknowledgements.}
The author thanks the IberLEF 2026 MSLG-SPA task organisers for releasing
the dataset and evaluation infrastructure.

% ---------------------------------------------------------------
\bibliographystyle{plain}
\begin{thebibliography}{99}

\bibitem{laraortiz2025}
Lara-Ortiz, M., et al.\ (2025).
\emph{A Corpus and Evaluation Framework for Mexican Sign Language Gloss
Translation}.
Proceedings of IberLEF 2026, CEUR Workshop Proceedings.

\bibitem{liu2020mbart}
Liu, Y., Gu, J., Goyal, N., Li, X., Edunov, S., Ghazvininejad, M.,
Lewis, M., \& Zettlemoyer, L.\ (2020).
\emph{Multilingual Denoising Pre-training for Neural Machine Translation}.
Transactions of the Association for Computational Linguistics, 8, 726--742.

\bibitem{hu2022lora}
Hu, E.\ J., Shen, Y., Wallis, P., Allen-Zhu, Z., Li, Y., Wang, S.,
Wang, L., \& Chen, W.\ (2022).
\emph{LoRA: Low-Rank Adaptation of Large Language Models}.
ICLR 2022.

\bibitem{sennrich2016bt}
Sennrich, R., Haddow, B., \& Birch, A.\ (2016).
\emph{Improving Neural Machine Translation Models with Monolingual Data}.
ACL 2016, pp.\ 86--96.

\bibitem{wslp2025}
Anonymous.\ (2025).
\emph{Fine-tuning Pre-trained Language Models for Bidirectional Sign
Language Gloss Translation}.
Proceedings of WSLP 2025.
% Replace with actual citation when available.

\bibitem{hoang2018iterative}
Hoang, V.\ C.\ D., Koehn, P., Haffari, G., \& Cohn, T.\ (2018).
\emph{Iterative Back-Translation for Neural Machine Translation}.
Proceedings of the 2nd Workshop on Neural Machine Translation and
Generation, pp.\ 18--24.

\end{thebibliography}

\end{document}
